\section{Experimental Setup} \label{sec:experimental}

\subsection{Experimental Design}

We conduct a systematic comparative evaluation using controlled experimental conditions to ensure fair assessment:

\subsubsection{Controlled Variables}

\begin{enumerate}
  \item \textbf{Data source} --- Identical dataset across all three models
  \item \textbf{Train/val/test splits} --- Same stratified partition for reproducibility
  \item \textbf{Preprocessing} --- Identical preprocessing pipeline for all approaches
  \item \textbf{Evaluation metrics} --- Computed using identical methodology
  \item \textbf{Random seeds} --- Fixed to eliminate stochastic variation
\end{enumerate}

\subsubsection{Independent Variables}

\begin{enumerate}
  \item \textbf{Approach type} --- TF-IDF (classical ML), BiLSTM (RNN), BERT (Transformer)
  \item \textbf{Model architecture} --- Varies per approach (vectorizer vs neural network vs pre-trained)
  \item \textbf{Training paradigm} --- Supervised fine-tuning vs transfer learning
\end{enumerate}

\subsubsection{Dependent Variables}

\begin{enumerate}
  \item \textbf{Accuracy, Precision, Recall, F1, AUC} --- Performance metrics
  \item \textbf{Inference latency} --- Model response time per prediction
  \item \textbf{Memory footprint} --- Model size and GPU memory requirements
  \item \textbf{Training time} --- Duration from initialization to convergence
\end{enumerate}

\subsection{Validation Strategy}

\subsubsection{Temporal Validation}

Models evaluated on temporally-held-out test set (15\% data, time-ordered split) ensuring temporal consistency: models not trained on future data.

\subsubsection{Cross-Validation}

For robustness analysis, stratified 5-fold cross-validation performed on full dataset with reported metrics as mean $\pm$ standard deviation across folds.

\subsubsection{Statistical Testing}

Paired t-tests (two-tailed, $\alpha = 0.05$) assess statistical significance of performance differences between approaches.

\subsection{Model Convergence Criteria}

\subsubsection{TF-IDF}

No training phase; convergence immediate upon vectorizer initialization. Validation on held-out validation set to tune threshold parameter.

\subsubsection{BiLSTM}

Training stops when:
\begin{enumerate}
  \item Validation accuracy plateaus for 2 consecutive epochs, OR
  \item Maximum 10 epochs reached
\end{enumerate}

Learning curves (training/validation loss) monitored for overfitting detection.

\subsubsection{BERT}

Fine-tuning stops when:
\begin{enumerate}
  \item Validation accuracy plateaus for 1 epoch, OR
  \item Maximum 5 epochs reached
\end{enumerate}

Early stopping checkpoint loaded containing weights with best validation performance.

\subsection{Threshold Optimization}

For binary classification, optimal decision thresholds determined via:

\begin{enumerate}
  \item \textbf{ROC curve analysis} --- Maximize Youden's index: $J = \text{Sensitivity} + \text{Specificity} - 1$
  
  \item \textbf{F1 optimization} --- Select threshold maximizing F1-Score on validation set
  
  \item \textbf{Precision-recall trade-off} --- Institution-specific tuning based on acceptable false positive rate
\end{enumerate}

Thresholds determined on validation set and applied unchanged to test set for unbiased evaluation.

\subsection{Generalization Analysis}

\subsubsection{In-Domain Testing}

Models evaluated on academic text from same domain (AI/ML papers) used in training.

\subsubsection{Out-of-Domain Testing}

Model robustness evaluated on:
\begin{itemize}
  \item Different academic domains (biology, chemistry papers)
  \item Non-academic technical writing (documentation, blog posts)
  \item Adversarial plagiarism (extracted from plagiarism detection literature)
\end{itemize}

\subsubsection{Plagiarism Type Breakdown}

Performance analyzed separately for each plagiarism category:
\begin{itemize}
  \item Direct copying detection rate
  \item Paraphrasing detection rate
  \item Semantic variation detection rate
\end{itemize}

This breakdown identifies model strengths/weaknesses for specific plagiarism types.

\subsection{Ablation Studies}

For transformer model, ablation analysis examines:

\begin{enumerate}
  \item \textbf{Pre-training value} --- Fine-tuning vs training from random initialization
  
  \item \textbf{Model size} --- bert-base vs bert-large comparison
  
  \item \textbf{Task head complexity} --- Single FC layer vs multi-layer MLP
  
  \item \textbf{Sequence length} --- Impact of truncation to 128 tokens
\end{enumerate}

Results quantify contribution of each component to final performance.

\subsection{Error Analysis}

Detailed analysis of misclassifications:

\begin{enumerate}
  \item \textbf{False positives} --- Legitimate original content flagged as plagiarism (unnecessary accusations)
  
  \item \textbf{False negatives} --- Undetected plagiarism (missed violations)
  
  \item \textbf{Confidence distribution} --- Histogram of model confidence scores for correct vs incorrect predictions
  
  \item \textbf{Challenging cases} --- Identify dataset instances where multiple models struggle
\end{enumerate}

This analysis guides threshold setting and future improvement directions.
